\begin{center}
\resizebox{0.96\linewidth}{!}{%
\begin{tikzpicture}[font=\small,>=Latex]
  \path[draw=GrisOscuro,very thick,rounded corners=8pt,fill=Gris]
    (0.20,0.35) rectangle (14.70,5.45);
  \node[anchor=north west,font=\bfseries] at (0.45,5.25)
    {espacio funcional \(\mathfrak C\) (representación esquemática)};

  % Cuenca funcional: la lengua superior entra en U_e sin tocar el corte cercano.
  \path[draw=Verde,very thick,fill=VerdeClaro]
    (4.35,3.05)
    .. controls (5.05,3.85) and (6.15,3.90) .. (7.05,3.30)
    .. controls (8.15,4.70) and (11.85,4.80) .. (13.55,3.65)
    .. controls (14.15,3.20) and (14.05,1.70) .. (13.20,1.25)
    .. controls (11.90,0.65) and (9.10,0.90) .. (7.95,1.75)
    .. controls (7.25,2.25) and (7.00,2.65) .. (6.35,2.75)
    .. controls (5.55,2.88) and (4.85,2.55) .. (4.35,3.05) -- cycle;
  \node[Verde,align=center] at (10.70,3.55)
    {cuenca funcional\\\(\mathcal B_{\mathfrak C}(A)\)};
  \node[draw=Verde,very thick,circle,fill=white,minimum size=8mm] (A) at (12.35,2.55) {\(A\)};

  % Corte finito-dimensional de perfiles constantes.
  \draw[GrisOscuro,very thick]
    (0.85,1.55) .. controls (4.15,1.95) and (8.40,1.35) .. (13.95,1.75);
  \node[anchor=west,fill=white,inner sep=2pt] (slice) at (0.70,2.55)
    {sección de reinicio \(\iota(\mathbb R^n)\)};
  \draw[-{Latex[length=1.7mm]},GrisOscuro] (slice.south)--(1.65,1.66);
  \draw[Verde,line width=3pt]
    (8.30,1.50) .. controls (10.10,1.40) and (12.05,1.48) .. (13.55,1.69);
  \node[Verde,anchor=south,fill=white,inner sep=1pt] at (10.65,1.62)
    {\(\iota[\mathcal B_{\mathrm{reset}}(A)]\)};

  % Equilibrio y vecindad funcional.
  \node[draw=Azul,very thick,circle,fill=AzulClaro,minimum size=7mm] (Ee) at (3.70,1.82) {\(E_e\)};
  \draw[Azul,very thick,dashed] (3.70,1.82) ellipse (1.25cm and 1.42cm);
  \node[Azul,anchor=east,align=right] at (2.35,3.00)
    {vecindad funcional\\\(U_e\)};
  \draw[-{Latex[length=2mm]},Azul] (2.45,2.83)--(2.85,2.55);
  \node[Naranja,align=center] at (5.15,4.40)
    {perfil no constante cercano\\que pertenece a la cuenca};
  \draw[-{Latex[length=2mm]},Naranja,thick] (5.15,4.05)--(4.72,3.12);

  \node[draw=GrisOscuro,rounded corners=2pt,fill=white,align=center] at (7.45,0.60)
    {ocultedad funcional \(\Longrightarrow\) ocultedad de reinicio;
     el recíproco no vale en general};
\end{tikzpicture}%
}
\par\smallskip
\begin{minipage}{0.92\linewidth}
\small\textbf{Esquema, no cuenca calculada.}
La figura muestra una configuración compatible con la jerarquía lógica: una
bola del corte de reinicio puede evitar la cuenca mientras una perturbación
funcional cercana sí la alcanza. No atribuye esta geometría a PLL, MAVPD, Wu ni
a otro caso del banco actual.
\end{minipage}
\end{center}
